\chapter{Acute Kidney Injury}

\section*{Definition}
Acute kidney injury (AKI) is a sudden -- usually over hours to days -- drop in how well your kidneys filter waste from your blood. The loss of filtration can be mild and silent or severe enough to require dialysis. It is diagnosed by a rise in creatinine (a waste product of muscle) and/or a fall in urine output. The RIFLE and KDIGO systems grade severity from "risk/injury" to "failure," helping doctors predict who needs closer monitoring.

\section*{When to See a Doctor}
Seek urgent care if:
\begin{itemize}
  \item You are passing little or no urine
  \item You have swelling in your legs, ankles, or face, or sudden breathlessness
  \item You feel confused, drowsy, or have a fluttering or racing heartbeat (possible high potassium)
  \item You are vomiting persistently and cannot keep fluids down
  \item You see blood in your urine or have severe flank pain
  \item You have recently started a new medicine, had an imaging study with contrast, or are severely dehydrated and feel unwell
\end{itemize}

\section*{How It Develops}
Your kidneys filter about 180 liters of blood daily through millions of microscopic filters called nephrons. In AKI, blood flow to the kidneys falls, or the filters themselves are damaged, or urine flow out of the kidneys becomes blocked. When filtration drops, creatinine, urea, potassium, and phosphate accumulate in the blood while the kidneys lose their ability to balance fluid and electrolytes. Because the injured kidney can no longer concentrate urine normally, you may not make much urine, and fluid backs up into your lungs and legs. If potassium rises too high, the heart rhythm can become dangerous, which is why AKI is a medical emergency. The kidneys often recover if the underlying cause is corrected quickly, but severe injury can leave permanent damage or be fatal.

\section*{Causes}
\begin{itemize}
  \item Reduced blood flow to the kidneys: severe dehydration, bleeding, heart failure, low blood pressure, sepsis, major surgery, or burns (pre-renal AKI)
  \item Direct kidney damage: nephrotoxic medicines such as NSAIDs, certain antibiotics, and IV contrast dye; rhabdomyolysis; and glomerulonephritis (intrinsic AKI)
  \item Blockage of urine flow: enlarged prostate, kidney stones, or a tumor blocking both ureters or the bladder outlet (post-renal AKI)
  \item Chronic conditions: diabetes, high blood pressure, and existing kidney disease make the kidneys more fragile and AKI more likely
\end{itemize}

\section*{Related Symptoms}
\begin{itemize}
  \item Decreased urine output or no urine at all
  \item Swelling of the legs, ankles, or face
  \item Shortness of breath or chest pressure from fluid overload
  \item Fatigue, weakness, and confusion
  \item Nausea, vomiting, and loss of appetite
  \item Flank or back pain
  \item Muscle cramps (from electrolyte imbalance)
\end{itemize}

\section*{Medical Evaluation}
\begin{itemize}
  \item Blood tests for creatinine, urea, electrolytes (especially potassium), and blood gases to grade the severity of kidney failure.
  \item Urine tests, including microscopy and a urine sodium level, to distinguish low blood flow from direct kidney damage.
  \item A kidney ultrasound to rule out blockage of urine flow and assess kidney size and structure.
  \item A review of medications, recent procedures or contrast exposure, blood pressure, volume status, and underlying heart or kidney disease.
  \item In uncertain cases, a kidney biopsy may be needed, especially when glomerulonephritis or other intrinsic disease is suspected.
\end{itemize}

\section*{Treatment Options}
\begin{itemize}
  \item Intravenous fluids to restore blood flow in dehydration or blood loss, given carefully in people who cannot excrete urine well.
  \item Stopping nephrotoxic medicines and adjusting doses of drugs cleared by the kidneys.
  \item Treating the underlying cause: antibiotics for sepsis, blood products for bleeding, relieving a urinary obstruction with a catheter or stent, and treating heart failure.
  \item Medicines to manage potassium (insulin and glucose, calcium, potassium-binding drugs) and other electrolyte abnormalities.
  \item Dialysis (hemodialysis or continuous renal replacement) when fluid overload, very high potassium, or severe acid buildup threaten your life.
  \item Close follow-up after recovery, because even a reversible AKI increases the future risk of chronic kidney disease.
\end{itemize}


\section*{Self-Care and Home Management}
\begin{itemize}
  \item Avoid NSAIDs (ibuprofen, naproxen) when you are dehydrated or feeling unwell, since they reduce blood flow to the kidneys.
  \item Stay well hydrated with water during illness, heat, or diarrhea, and replace electrolytes if vomiting or sweating heavily.
  \item Do not stop prescribed blood-pressure or heart medicines on your own -- check with your doctor first.
  \item Monitor your urine output and weight; sudden weight gain can signal fluid retention.
  \item AKI is not manageable at home once red flags appear -- go to the emergency department rather than waiting.
\end{itemize}

\section*{References}
\begin{thebibliography}{99}
\bibitem{acute_kidney_injury:kdogi} Khwaja A. KDIGO clinical practice guidelines for acute kidney injury. \textit{Nephron Clin Pract}. 2012;120:c179--c184.
\bibitem{acute_kidney_injury:rifle} Bellomo R, Ronco C, Kellum JA, et al. Acute renal failure -- definition, outcome measures, animal models, fluid therapy and information technology needs. \textit{Crit Care}. 2004;8:R204--R212.
\bibitem{acute_kidney_injury:kdigo-summary} Kellum JA, Lameire N, Aspelin P, et al. Diagnosis, evaluation, and management of acute kidney injury: a KDIGO summary. \textit{Intensive Care Med}. 2013;39:1586--1598.
\bibitem{acute_kidney_injury:epidemiology} Hoste EAJ, Kellum JA, Selby NM, et al. Global epidemiology and outcomes of acute kidney injury. \textit{Nat Rev Nephrol}. 2018;14:607--625.
\end{thebibliography}
